\section{Prototype Implementation}
\label{sec:implementation}
Let us start with the general implementation of the application, where we have our src folder, this is the backend, and the poll-frontend, being respectively the frontend.

\subsection{Frontend}
The frontend is Vite and React, where it is a single-page application using React where the core domain logic is within App.jsx.
Using Reacts useState  we can keep track of Polls, votes, users and even users login status also anything else that needs to be kept track of that is dynamic.
Handles users interaction; 1) creation of users, 2) creation of polls and 3) voting. Creation of Polls is set to users that have the role of ADMIN but this is checked from the backend using the credentials used from here. How is it handled in the frontend? If the user is not of ADMIN [role], then it does not give them the button to create a poll (see image).
\lstinputlisting[language=html]{code/short_sample_of_app_jsx.html}

Also upon voting and creation of polls one obviously needs to be logged in, so this is handled.
Then it uses fetch in functions to make request for creating users, creating polls, voting and more.

\subsection{Backend}

The backend is a Spring Boot application that provides a REST API for the frontend. The structure is organized into several key packages:

\begin{itemize}
    \item Domain: This package contains the data models (also known as entities) like User, Poll, and Vote.  Using JPA we can save all these entities into the H2 database directly to a table in the database. You are provided with all different functions and methods, using the $@data$ annotation, thereby effectively interacting with the user database.
    \item Config, repository and cache: These are the folders mainly responsible for the caching Logic using neo4j.
    \item Controller: This package exposes the API endpoints. Classes like UserController, PollController, and VoteController handle incoming HTTP requests from the frontend. For example, when you log in or create a poll, these controllers process that request. Using the springs annotations like $@GetMapping$, $@RequestMapping$ and $@PostMapping$
    \item Security: This package manages user authentication and authorization using Spring Security. It ensures that users are who they say they are and controls access to different parts of the API. For instance, it restricts poll creation to users with an ADMIN role.
    \item Service: This package holds the application's business logic. It includes the PollManager which is responsible for our controller logic for Polls and integrating with messaging systems like RabbitMQ.
\end{itemize}



\subsection{Security}

Here we will go trough to our implementation of the security aspect of the project, starting off with our Role.java enum which is simple, and where a USER can have the Role of "NORMAL" or "ADMIN".

Now that we have this Role, we can use it in our User class that each User has a role, or just set it as a Set of Role(s).
Then our most important part (perhaps) is the webconfig, this in turns dictates what is required to access certain information.
\lstinputlisting[language=java]{code/WebSecurityConfig_shorter.java}
The code is quite easy to understand, as i have added quite useful comments to it, the most important takeaway, is that ADMIN roles are allowed to create and delete Polls and the rest both Normal and Admin users can do. 

Inside the User controller we have a /me path, to easily retrieve information, like username authentication (role) and password. 

Then to have things secure in our application we can easily just add the 	@PreAuthorize("hasRole('ADMIN')") for the POST and DELETE requests inside the PollController, the rest are after the requirements of the `WebSecurityConfig.java` file.
Now the ADMIN is created at the start of the application.
This way an ADMIN must be preconfigured to be made at startup, and in the command-line one can see the password and username set, which they can use to directly log in if wanted.




\subsection{Cache Implementation}

The implementation of the Neo4j cache required quite a lot of modification to the original code. We added several new domain classes (PollAggregate and OptionVote), a specific Neo4jPollCache class which is an implementation of an interface PollCache, and a PollCacheRepository, which extends the Neo4jRepository from Spring Data.

\subsubsection{PollAggregate}
The PollAggregate is a specific object designed to be used to cache the votes of the specific options of a poll in a single object. It is identified by the Poll it is created from and contains a list of OptionVote, which is explained more in section~\ref{sec:OptionVote}. It also contains a LocalDateTime object which is used to decide whether the poll should be removed from the cache or not.

\subsubsection{OptionVote}
\label{sec:OptionVote}
The OptionVote object contains the information we need to aggregate the number of votes for the options of a specific poll. It contains an option of the poll and the number of votes that the option has. 

\subsubsection{Neo4jPollCache}
This is the cache manager. It is connected to the PollCacheRepository, and to RabbitMQ. Whenever a user votes on a poll, this object checks whether there exists a PollAggregate object in the cache and whether this aggregate contains the option that is voted on. If everything matches, it either decrements or increments the results based on the user input. Then it updates the LocalDateTime object in PollAggregate and puts it back into the cache. If the PollAggregate does not exist, i.e. the poll has never been voted on or the poll has been cleared from the cache, it creates a new PollAggregate object corresponding to the vote and puts it into the cache. If the PollAggregate exists but the option that is voted on does not exist in the aggregate, it creates a new OptionVote object with the new vote and adds it to the existing PollAggregate, which is then put back into the cache.   

It also has a scheduled cleaning method that runs every minute, removing polls that haven't been interacted with through a method deleteAggregateWithOptions defined in PollCacheRepository.

\subsubsection{PollCacheRepository}
This is the actual cache. It extends Neo4jRepository with a method for removing the polls through a Cypher query.



EXAMPLE SCRIPT BELOW.
This section should provide brief details of how the prototype has been implemented.
You may want to use come code snippets here, but only focus on core features and aspects.
You are not meant to copy/paste your whole application code into the report.
Focus for instance how other developers may run your application and how they might develop it further...


The example below shows how you may include code. There are similar
styles for many other langages - in case you do not use Java in your
project. You can wrap the listing into a figure in case you need to
refer to it. How to create a figure was shown in Section~\ref{sec:technology}.

\lstinputlisting[language=java]{code/BoksVolum.java}
